%%%%%%%%%%%%%%%%%%%%%%%%%%%%%%%%%%%%%%%%%
% Author: Huseyin Ugur Yildiz
%%%%%%%%%%%%%%%%%%%%%%%%%%%%%%%%%%%%%%%%%
% GENERATED — do not edit main.tex by hand. Edit _data/cv.yml and
% _data/publications.yml, then regenerate with scripts/render_cv_tex.py.

\documentclass[letterpaper,11pt]{article}
\usepackage[empty]{fullpage}
\usepackage{titlesec}
\usepackage{marvosym}
\usepackage[usenames,dvipsnames]{color}
\usepackage{enumitem}
\usepackage{fancyhdr}
\usepackage[hidelinks]{hyperref}
% Document properties: without these the PDF reaches a jury or a funding panel with an
% empty Title/Author, so a reader's file manager and PDF viewer label it by filename.
\hypersetup{
  pdftitle={Curriculum Vitae --- Huseyin Ugur Yildiz},
  pdfauthor={Huseyin Ugur Yildiz},
  pdfsubject={Academic curriculum vitae: appointments, education, publications, teaching, supervision and service},
  pdfkeywords={Wireless Ad Hoc and Sensor Networks, Underwater Acoustic Sensor Networks, Internet of Things, Drone Networks, Smart Grids, Operations Research, Metaheuristics, Neural Networks, Reinforcement Learning},
  pdfcreator={pdflatex (generated from repository data, not hand-edited)},
  pdflang={en}
}
\usepackage{ragged2e}
\usepackage{bold-extra}
\usepackage[english]{babel}
\usepackage{fontawesome5}
\pagenumbering{arabic}
\usepackage{lastpage}
\usepackage{comment}
\usepackage[T1]{fontenc}
\usepackage[utf8]{inputenc}

%----------FONT OPTIONS----------
% sans-serif
% \usepackage[sfdefault]{FiraSans}
% \usepackage[sfdefault]{roboto}
% \usepackage[sfdefault]{noto-sans}
% \usepackage[default]{sourcesanspro}

% serif
% \usepackage{CormorantGaramond}
% \usepackage{charter}

\pagestyle{fancy}
\fancyhf{} % clear all header and footer fields
\fancyfoot[L]{\small \today}
\fancyfoot[C]{\small Huseyin Ugur Yildiz}
\fancyfoot[R]{\small \thepage/\pageref{LastPage}}
\setlength{\footskip}{14pt}
\renewcommand{\headrulewidth}{0pt}
\renewcommand{\footrulewidth}{0pt}

% Adjust margins
\addtolength{\oddsidemargin}{-0.5in}
\addtolength{\evensidemargin}{-0.5in}
\addtolength{\textwidth}{1in}
\addtolength{\topmargin}{-.6in}
\addtolength{\textheight}{1in}
\setlength{\headwidth}{\textwidth}

%\urlstyle{same}
\raggedbottom
\raggedright
\setlength{\tabcolsep}{0in}

% Sections formatting
\titleformat{\section}{
  \vspace{-4pt}\Large\bf\scshape 
}{}{0em}{}[\color{black}\titlerule \vspace{-5pt}]

% Ensure that generate pdf is machine readable/ATS parsable
\input{glyphtounicode}
\pdfgentounicode=1

%%%%%%%%%%%%%%%%%%%%%%%%%%%%%%%%%%%%%%%%%
%----------CV STARTS HERE------------
%%%%%%%%%%%%%%%%%%%%%%%%%%%%%%%%%%%%%%%%%
\begin{document}

\begin{center}
{\huge \scshape \textbf{Huseyin Ugur Yildiz, Ph.D.}} \\ \vspace{5pt}

{\scshape \small Associate Professor of Electrical Engineering | IEEE Senior Member | Operations Research \\ · Network Optimization · ML/RL | Wireless, Underwater \& Aerial Networks} \\ \vspace{5pt}


%{\scshape \small Operations Research Scientist | Expert in MILP, Heuristics, Network Flow | Ph.D. in Electrical Engineering | IEEE Senior Member
%} \\ \vspace{5pt}

%{\scshape \small Researcher with 10+ Years of Experience Applying Operations Research Techniques to Terrestrial and Underwater Acoustic Wireless Sensor Networks, IEEE Senior Member
%} \\ \vspace{5pt}

\small \raisebox{-0.1\height}\faPhone~+90 (312) 585 02 21~\textbar~
\href{mailto:hugur.yildiz@tedu.edu.tr}{\raisebox{-0.2\height}\faEnvelope~\underline{hugur.yildiz@tedu.edu.tr}} ~\textbar~
\href{https://www.linkedin.com/in/huguryildiz/}{\raisebox{-0.2\height}\faLinkedinIn~\underline{huguryildiz}} ~\textbar~
\href{https://www.huguryildiz.com/}{\raisebox{-0.2\height}\faGlobeAmericas~\underline{huguryildiz.com}} ~\textbar~
\href{https://goo.gl/maps/vU69BadSTiav3s3B9}{\raisebox{-0.2\height}{\faMapMarker*}~Ankara, Türkiye}
%\faOrcid~\href{https://orcid.org/0000-0002-1556-2634}{0000-0002-1556-2634} ~\textbar~
%\faResearchgate~\href{https://www.researchgate.net/profile/Huseyin-Yildiz-3}{ResearchGate}

\end{center}
\vspace{-5pt}

%%%%%%%%%%%%%%%%%%%%%%%%%%%%%%%%%%%%%%%%%
%-----------SUMMARY--------------------
%%%%%%%%%%%%%%%%%%%%%%%%%%%%%%%%%%%%%%%%%
\vspace{-10pt}
\section{\faUser~Summary}
\justify
Associate Professor of Electrical and Electronics Engineering with a Ph.D. specializing in operations research and optimization for wireless, underwater, and aerial communication systems. Over 10 years of research experience, I develop mathematical optimization models for routing, resource allocation, reliability, and energy efficiency across complex networked systems, integrating reinforcement learning techniques for dynamic and large-scale environments. My current interests include hybrid optimization–learning methodologies for next-generation network paradigms, including quantum communication systems. I have authored 40+ peer-reviewed journal and conference publications and actively serve as editor, reviewer, and TPC member. I also teach probability, signals and systems, and communication networks, contributing to curriculum development, accreditation, and academic leadership.
\begin{comment}
\begin{itemize}[leftmargin=*]
\setlength\itemsep{-0.4em}
\item \noindent Ph.D. in Electrical and Electronics Engineering, specializing in operations research for wireless sensor networks.
\item \noindent 10+ years of experience in optimizing terrestrial and underwater acoustic wireless sensor networks.
\item \noindent 6+ years of teaching experience in courses such as signal processing, probability, and communication systems.
\item \noindent 5+ years of industry experience in network design, operations, and security.
\item \noindent Author of 20+ journal articles (14 in IEEE), 10+ conference papers, with over 20 first-author contributions. 
\item \noindent IEEE Senior Member ($<$ 10\% of the total IEEE membership achieves this grade).
%\item \noindent Experience with presenting, briefing, or communicating analytical research material in industry or academia.
\end{itemize}

\end{comment}

%%%%%%%%%%%%%%%%%%%%%%%%%%%%%%%%%%%%%%%%%
%-----------EXPERIENCE-------------------
%%%%%%%%%%%%%%%%%%%%%%%%%%%%%%%%%%%%%%%%%
\vspace{-10pt}
\section{\faIndustry~Professional Experience}
\justify

\noindent \textbf{\href{https://www.tedu.edu.tr/en}{TED University}} \hfill {Ankara, Türkiye} \\
\noindent {Associate Professor of Electrical and Electronics Engineering} \hfill {03/2021 -- Present} \\
\noindent {Chair of the Electrical and Electronics Engineering Department} \hfill {07/2021 -- 07/2024} \\
\noindent {Assistant Professor of Electrical and Electronics Engineering} \hfill {09/2016 -- 03/2021} \\
\vspace{-20pt}
\begin{itemize}[leftmargin=*]
\setlength\itemsep{-0.5em}
\item \noindent Conducted research on network flow–based optimization to enhance wireless sensor network performance; authored 40+ peer-reviewed publications, including 15 articles in IEEE journals.
\item \noindent Led academic operations for 100+ students, 8 faculty members, and 4 research assistants; oversaw accreditation and curriculum, established a new teaching laboratory, launched an M.Sc. program, and founded the IEEE Student Branch.
\item \noindent Developed \href{https://vera-eval.app/}{\underline{VERA}}, a web-based capstone jury evaluation platform for secure rubric-based assessment, analytics, and reporting, and \href{https://github.com/huguryildiz/KAIROS}{\underline{KAIROS}}, an OR-Tools CP-SAT university timetabling system with repair search and a separate post-solve validator.
\item \noindent Taught undergraduate and graduate courses in communications, signal processing, and probability using active learning methodologies; mentored 40+ senior design teams and supervised 8 graduate theses.
\end{itemize}

\noindent \textbf{\href{https://www.tusas.com/en}{Turkish Aerospace}} \hfill {Ankara, Türkiye} \\
\noindent {Avionics Design Engineer, Unmanned Aerial Vehicles} \hfill {12/2015 -- 08/2016} \\
\vspace{-20pt}
\begin{itemize}[leftmargin=*]
\setlength\itemsep{-0.5em}
\item \noindent Involved in development and integration of network architectures for the \href{https://www.tusas.com/en/products/uav/operative-strategic-uav-systems/anka}{\underline{ANKA}} UAV program; built and configured UAV network infrastructure and designed high-availability Cisco systems ensuring 24/7 readiness.
\end{itemize}

\noindent \textbf{\href{https://www.turktelekom.com.tr/en}{Turk Telekom}} \hfill {Ankara, Türkiye} \\
\noindent {Senior Network Engineer, Deep Packet Inspection} \hfill {12/2010 -- 11/2015} \\
\vspace{-20pt}
\begin{itemize}[leftmargin=*]
\setlength\itemsep{-0.5em}
\item \noindent Contributed to design and operation of the \href{https://www.guvenlinet.org.tr/}{\underline{Safer Internet Service}}, a nationwide content-filtering platform serving 5M+ users; managed 100+ Cisco/Procera DPI systems, ensuring 99.9\% uptime and regulatory compliance.
\end{itemize}

%%%%%%%%%%%%%%%%%%%%%%%%%%%%%%%%%%%%%%%%%
%-----------EDUCATION--------------------
%%%%%%%%%%%%%%%%%%%%%%%%%%%%%%%%%%%%%%%%%
\vspace{-10pt}
\section{\faGraduationCap~Education}
\justify

\noindent \textbf{\href{https://www.etu.edu.tr/en}{TOBB University of Economics and Technology (ETU)}} \hfill {Ankara, Türkiye} \\
\noindent {Ph.D., Electrical and Electronics Engineering} \hfill  {01/2014 -- 04/2016} \\
\noindent \emph{Dissertation Title:} \href{https://doi.org/10.1109/JSEN.2015.2486960}{\underline{Transmission power control for link–level handshaking in wireless sensor networks}} \\
%\noindent \emph{Doctoral Advisor:} Prof. Bulent Tavli\\

\vspace{-10pt}

\noindent \textbf{\href{https://www.etu.edu.tr/en}{TOBB University of Economics and Technology}} \hfill {Ankara, Türkiye} \\
\noindent {M.Sc., Electrical and Electronics Engineering} \hfill {09/2011 -- 09/2013} \\
\noindent \emph{Thesis Title:} \href{https://doi.org/10.1016/j.adhoc.2015.08.026}{\underline{Commun./comput. trade-offs in WSNs: Comparing network-level and node-level strategies}} \\
%\noindent \emph{Advisor:} Prof. Bulent Tavli, \emph{Co-Advisor:} Prof. Kemal Bicakci \\

\vspace{-10pt}

\noindent \textbf{\href{https://w3.bilkent.edu.tr/bilkent/}{Bilkent University}} \hfill {Ankara, Türkiye} \\
\noindent {B.Sc., Electrical and Electronics Engineering} \hfill {08/2005 -- 08/2009}

%%%%%%%%%%%%%%%%%%%%%%%%%%%%%%%%%%%%%%%%%
%-----------TECHNICAL SKILLS-------------------
%%%%%%%%%%%%%%%%%%%%%%%%%%%%%%%%%%%%%%%%%
\vspace{-10pt}
\section{\faCogs~Technical Skills}
\justify

\begin{itemize}[leftmargin=*]
\setlength\itemsep{-0.5em}
\item \noindent \textbf{Optimization \& Operations Research:} 
Linear Programming (LP), Mixed-Integer Programming (MIP), Multi-Objective Optimization, Network Flow Programming, Constraint Programming, Metaheuristics; GAMS, Gurobi, CPLEX, XPRESS, PuLP, Pyomo.
\item \noindent \textbf{Machine Learning \& Reinforcement Learning:} 
Supervised and Unsupervised Learning, Regression and Classification Models, Decision Trees, Neural Networks, Markov Decision Processes (MDPs), Q-Learning, Proximal Policy Optimization (PPO), Deep Reinforcement Learning.
\item \noindent \textbf{Programming \& Scientific Computing:} 
Python (NumPy, Pandas, Scikit-Learn, Matplotlib), MATLAB-Octave, Simulink, \LaTeX.
\end{itemize}

%%%%%%%%%%%%%%%%%%%%%%%%%%%%%%%%%%%%%%%%%
%-----------SELECTED RESEARCH SOFTWARE---
%%%%%%%%%%%%%%%%%%%%%%%%%%%%%%%%%%%%%%%%%
\vspace{-10pt}
\section{\faCodeBranch~Selected Research Software}
\justify

\begin{itemize}[leftmargin=*]
\setlength\itemsep{-0.5em}
\item \noindent \textbf{\href{https://github.com/huguryildiz/research-graph}{\underline{research-graph}}} --- MIT-licensed Python verifier for multi-agent research: schema validation, provenance-chain and producer--reviewer role checks, and bounded revisions; public beta with automated tests and continuous integration.
\item \noindent \textbf{\href{https://github.com/huguryildiz/KAIROS}{\underline{KAIROS}}} --- OR-Tools CP-SAT university timetabling system with repair search, independent post-solve constraint validation, a bilingual web interface, and structured schedule exports.
\end{itemize}
\vspace{-5pt}

%%%%%%%%%%%%%%%%%%%%%%%%%%%%%%%%%%%%%%%%%
%-----------CERTIFICATES-------------------
%%%%%%%%%%%%%%%%%%%%%%%%%%%%%%%%%%%%%%%%%
\vspace{-10pt}
\section{\faTrophy~Certificates}
\justify

\begin{itemize}[leftmargin=*]
\setlength\itemsep{-0.5em}
\item \noindent \textbf{Reinforcement Learning Specialization}, University of Alberta -- Coursera, 2025 (\href{https://www.coursera.org/account/accomplishments/specialization/26BOL8LTVD7S}{\underline{Credential}})
\item \noindent \textbf{Deep Learning Specialization}, DeepLearning.AI -- Coursera, 2025 (\href{https://www.coursera.org/account/accomplishments/specialization/JGGZ4NGSKTQZ}{\underline{Credential}})
\item \noindent \textbf{Machine Learning Specialization}, DeepLearning.AI -- Coursera, 2025 (\href{https://www.coursera.org/account/accomplishments/specialization/FGT2PNQ9NYVL}{\underline{Credential}})
\end{itemize}

%%%%%%%%%%%%%%%%%%%%%%%%%%%%%%%%%%%%%%%%%
%-----------AWARDS-------------------
%%%%%%%%%%%%%%%%%%%%%%%%%%%%%%%%%%%%%%%%%
\vspace{-10pt}
\section{\faMedal~Honors \& Awards}
\justify

\noindent [A2] \textbf{Elevation to the IEEE Senior Member grade,} an honor recognizing significant contributions to the profession (achieved by less than 10\% of IEEE members), 2021.
\vspace{5pt}

\noindent [A1] \textbf{The best paper award (3rd place),} “Utilization of multi-sink architectures for lifetime maximization in underwater sensor networks,” 2nd IEEE Middle East \& North Africa Communications Conference (MENACOMM'19), Manama, Bahrain, 2019.

%%%%%%%%%%%%%%%%%%%%%%%%%%%%%%%%%%%%%%%%%
%-----------RESEARCH INTERESTS-----------
%%%%%%%%%%%%%%%%%%%%%%%%%%%%%%%%%%%%%%%%%
\vspace{-10pt}
\section{\faMicroscope~Research Areas}
\justify

Wireless Ad Hoc and Sensor Networks, Underwater Acoustic Sensor Networks, Internet of Things, Drone Networks, Smart Grids, Operations Research, Metaheuristics, Neural Networks, Reinforcement Learning.

%%%%%%%%%%%%%%%%%%%%%%%%%%%%%%%%%%%%%%%%%
%-----------PUBLICATIONS-------------------
%%%%%%%%%%%%%%%%%%%%%%%%%%%%%%%%%%%%%%%%%
\vspace{-10pt}
\section{\faBookOpen~Publications}
\justify

\subsection*{(1) \underline{Journal Articles:}}

\noindent [J24] Tantur Karagul, C., Akgun, M. B., \textbf{Yildiz, H. U.}, \& Tavli, B. (2025). Mitigating energy cost of connection reliability in UWSNs through non-uniform k-connectivity. \emph{IEEE Internet of Things Journal, 12\nocorr}(22), 47817–47826.\par\vspace{3pt}\noindent [J23] Asci, M., Akusta Dagdeviren, Z., Akram, V. K., \textbf{Yildiz, H. U.}, Dagdeviren, O., \& Tavli, B. (2025). Enhancing drone network resilience: Investigating strategies for k-connectivity restoration. \emph{Computer Standards \& Interfaces, 92\nocorr}, 103941.\par\vspace{3pt}\noindent [J22] Gultekin, B., Nurcan-Atceken, D., Altin-Kayhan, A., \textbf{Yildiz, H. U.}, \& Tavli, B. (2023). Exploring the tradeoff between energy dissipation, delay, and the number of backbones for broadcasting in wireless sensor networks through goal programming. \emph{Ad Hoc Networks, 149\nocorr}, 103223.\par\vspace{3pt}\noindent [J21] \textbf{Yildiz, H. U.} (2023). Joint effects of void region size and sink architecture on underwater WSNs lifetime. \emph{IEEE Sensors Journal, 23\nocorr}(10), 11046–11056.\par\vspace{3pt}\noindent [J20] Cobanlar, M., \textbf{Yildiz, H. U.}, Akram, V. K., Dagdeviren, O., \& Tavli, B. (2022). On the tradeoff between network lifetime and k-connectivity-based reliability in UWSNs. \emph{IEEE Internet of Things Journal, 9\nocorr}(23), 24444–24452.\par\vspace{3pt}\noindent [J19] Carsancakli, M. F., Al Imran, M. A., \textbf{Yildiz, H. U.}, Kara, A., \& Tavli, B. (2022). Reliability of linear WSNs: A complementary overview and analysis of impact of cascaded failures on network lifetime. \emph{Ad Hoc Networks, 131\nocorr}, 102839.\par\vspace{3pt}\noindent [J18] Tekin, N., \textbf{Yildiz, H. U.}, \& Gungor, V. C. (2021). Node-level error control strategies for prolonging the lifetime of wireless sensor networks. \emph{IEEE Sensors Journal, 21\nocorr}(13), 15386–15397.\par\vspace{3pt}\noindent [J17] \textbf{Yildiz, H. U.}, Kurt, S., \& Tavli, B. (2019). Comparative analysis of transmission power level and packet size optimization strategies for WSNs. \emph{IEEE Systems Journal, 13\nocorr}(3), 2264–2274.\par\vspace{3pt}\noindent [J16] \textbf{Yildiz, H. U.} (2019). Maximization of underwater sensor networks lifetime via fountain codes. \emph{IEEE Transactions on Industrial Informatics, 15\nocorr}(8), 4602–4613.\par\vspace{3pt}\noindent [J15] Akbas, A., \textbf{Yildiz, H. U.}, Ozbayoglu, A. M., \& Tavli, B. (2019). Neural network–based instant parameter prediction for wireless sensor network optimization models. \emph{Wireless Networks, 25\nocorr}(6), 3405–3418.\par\vspace{3pt}\noindent [J14] Erdem, H. E., \textbf{Yildiz, H. U.}, \& Gungor, V. C. (2019). On the lifetime of compressive sensing based energy harvesting in underwater sensor networks. \emph{IEEE Sensors Journal, 19\nocorr}(12), 4680–4687.\par\vspace{3pt}\noindent [J13] Sayit, M., Cetinkaya, C., \textbf{Yildiz, H. U.}, \& Tavli, B. (2019). DASH–QoS: A scalable network layer service differentiation architecture for DASH over SDN. \emph{Computer Networks, 154\nocorr}, 12–25.\par\vspace{3pt}\noindent [J12] \textbf{Yildiz, H. U.} (2019). Investigation of maximum lifetime and minimum delay trade-off in underwater sensor networks. \emph{International Journal of Communication Systems, 32\nocorr}(7), e3924.\par\vspace{3pt}\noindent [J11] \textbf{Yildiz, H. U.}, Gungor, V. C., \& Tavli, B. (2019). Packet size optimization for lifetime maximization in underwater acoustic sensor networks. \emph{IEEE Transactions on Industrial Informatics, 15\nocorr}(2), 719–729.\par\vspace{3pt}\noindent [J10] \textbf{Yildiz, H. U.} (2018). The impact of transmission power levels set size on lifetime of wireless sensor networks in smart grids. \emph{Turkish Journal of Electrical Engineering \& Computer Sciences, 26\nocorr}(6), 3057–3071.\par\vspace{3pt}\noindent [J9] Yigit, M., \textbf{Yildiz, H. U.}, Kurt, S., Tavli, B., \& Gungor, V. C. (2018). A survey on packet size optimization for terrestrial, underwater, underground, and body area sensor networks. \emph{International Journal of Communication Systems, 31\nocorr}(11), e3572.\par\vspace{3pt}\noindent [J8] \textbf{Yildiz, H. U.}, Ciftler, B. S., Tavli, B., Bicakci, K., \& Incebacak, D. (2018). The impact of incomplete secure connectivity on the lifetime of wireless sensor networks. \emph{IEEE Systems Journal, 12\nocorr}(1), 1042–1046.\par\vspace{3pt}\noindent [J7] \textbf{Yildiz, H. U.}, Tavli, B., Kahjogh, B., \& Dogdu, E. (2017). The impact of incapacitation of multiple critical sensor nodes on wireless sensor network lifetime. \emph{IEEE Wireless Communications Letters, 6\nocorr}(3), 306–309.\par\vspace{3pt}\noindent [J6] Kurt, S., \textbf{Yildiz, H. U.}, Yigit, M., Tavli, B., \& Gungor, V. C. (2017). Packet size optimization in wireless sensor networks for smart grid applications. \emph{IEEE Transactions on Industrial Electronics, 64\nocorr}(3), 2392–2401.\par\vspace{3pt}\noindent [J5] Akbas, A., \textbf{Yildiz, H. U.}, Tavli, B., \& Uludag, S. (2016). Joint optimization of transmission power level and packet size for WSN lifetime maximization. \emph{IEEE Sensors Journal, 16\nocorr}(12), 5084–5094.\par\vspace{3pt}\noindent [J4] \textbf{Yildiz, H. U.}, Bicakci, K., Tavli, B., Gultekin, H., \& Incebacak, D. (2016). Maximizing wireless sensor network lifetime by communication/computation energy optimization of non-repudiation security service: Node level versus network level strategies. \emph{Ad Hoc Networks, 37\nocorr}(2), 301–323.\par\vspace{3pt}\noindent [J3] \textbf{Yildiz, H. U.}, Tavli, B., \& Yanikomeroglu, H. (2016). Transmission power control for link–level handshaking in wireless sensor networks. \emph{IEEE Sensors Journal, 16\nocorr}(2), 561–576.\par\vspace{3pt}\noindent [J2] \textbf{Yildiz, H. U.}, Temiz, M., \& Tavli, B. (2015). Impact of limiting hop count on the lifetime of wireless sensor networks. \emph{IEEE Communications Letters, 19\nocorr}(4), 569–572.\par\vspace{3pt}\noindent [J1] Batmaz, A. U., \textbf{Yildiz, H. U.}, \& Tavli, B. (2014). Role of unidirectionality and reverse path length on wireless sensor network lifetime. \emph{IEEE Sensors Journal, 14\nocorr}(11), 3971–3982.\par\vspace{-10pt}
\subsection*{(2) \underline{Editorial:}}

\noindent [E1] Haytaoglu, E., Arslan, S. S., Dagdeviren, O., \textbf{Yildiz, H. U.}, \& Ozturk, Y. (2025). Editorial brief for special issue: Mass connectivity and/or communication paradigms for the internet of things. \emph{Internet of Things, 32\nocorr}, 101625.
\vspace{-10pt}
\subsection*{(3) \underline{Conference Proceedings:}}

\noindent [C13] Tantur Karagul, C., Akgun, M. B., \textbf{Yildiz, H. U.}, \& Tavli, B. (2025, November). Non-uniform k-connectivity for energy-efficient and reliable underwater wireless sensor networks. In \emph{2025 33rd Telecommunications Forum (TELFOR)} (pp. 1–4). IEEE.\par\vspace{3pt}\noindent [C12] Un, B. E., \textbf{Yildiz, H. U.}, \& Tavli, B. (2021, May). Impact of critical node failures on lifetime of UWSNs with incomplete secure connectivity. In \emph{2021 IEEE International Black Sea Conference on Communications and Networking (BlackSeaCom)} (pp. 1–6). IEEE.\par\vspace{3pt}\noindent [C11] Ozmen, A., \textbf{Yildiz, H. U.}, \& Tavli, B. (2020, November). Impact of minimizing the eavesdropping risks on lifetime of underwater acoustic sensor networks. In \emph{2020 28th Telecommunications Forum (TELFOR)} (pp. 1–4). IEEE.\par\vspace{3pt}\noindent [C10] \textbf{Yildiz, H. U.} (2019, November). Utilization of multi-sink architectures for lifetime maximization in underwater sensor networks. In \emph{2019 2nd IEEE Middle East and North Africa COMMunications Conference (MENACOMM)} (pp. 1–5). IEEE.\par\vspace{3pt}\noindent [C9] \textbf{Yildiz, H. U.} (2019, October). Prolonging the lifetime of underwater sensor networks under sinkhole attacks. In \emph{14th International Conference on Underwater Networks \& Systems (WUWNet '19)} (pp. 1–5). ACM.\par\vspace{3pt}\noindent [C8] \textbf{Yildiz, H. U.}, Gungor, V. C., \& Tavli, B. (2018, June). A hybrid energy harvesting framework for energy efficiency in wireless sensor networks based smart grid applications. In \emph{2018 17th Annual Mediterranean Ad Hoc Networking Workshop (Med-Hoc-Net)} (pp. 1–6). IEEE.\par\vspace{3pt}\noindent [C7] Dagdeviren, O., Akram, V. K., Tavli, B., \textbf{Yildiz, H. U.}, \& Atilgan, C. (2016, October). Distributed detection of critical nodes in wireless sensor networks using connected dominating set. In \emph{2016 IEEE SENSORS} (pp. 1–3). IEEE.\par\vspace{3pt}\noindent [C6] Tantur, C., \textbf{Yildiz, H. U.}, Kurt, S., \& Tavli, B. (2016, October). Optimal transmission power level sets for lifetime maximization in wireless sensor networks. In \emph{2016 IEEE SENSORS} (pp. 1–3). IEEE.\par\vspace{3pt}\noindent [C5] \textbf{Yildiz, H. U.}, \& Tavli, B. (2015, December). Prolonging wireless sensor network lifetime by optimal utilization of compressive sensing. In \emph{2015 IEEE Globecom Workshops (GC Wkshps)} (pp. 1–6). IEEE.\par\vspace{3pt}\noindent [C4] \textbf{Yildiz, H. U.}, \& Tavli, B. (2014, December). The impact of random power assignment in handshaking on wireless sensor network lifetime. In \emph{2014 IEEE Globecom Workshops (GC Wkshps)} (pp. 201–206). IEEE.\par\vspace{3pt}\noindent [C3] \textbf{Yildiz, H. U.}, Kurt, S., \& Tavli, B. (2014, October). The impact of near-ground path loss modeling on wireless sensor network lifetime. In \emph{2014 IEEE Military Communications Conference (MILCOM)} (pp. 1114–1119). IEEE.\par\vspace{3pt}\noindent [C2] Akbas, A., \textbf{Yildiz, H. U.}, \& Tavli, B. (2014, May). Data packet length optimization for wireless sensor network lifetime maximization. In \emph{2014 10th International Conference on Communications (COMM)} (pp. 1–6). IEEE.\par\vspace{3pt}\noindent [C1] \textbf{Yildiz, H. U.}, Bicakci, K., \& Tavli, B. (2014, January). Communication/computation trade-offs in wireless sensor networks: Comparing network-level and node-level strategies. In \emph{2014 IEEE Topical Conference on Wireless Sensors and Sensor Networks (WiSNet)} (pp. 49–51). IEEE.\par\vspace{-10pt}
\subsection*{(4) \underline{Conference Proceedings (in Turkish):}}

\noindent [CT5] \textbf{Yildiz, H. U.} (2019, April). Improvement of underwater acoustic sensor networks performance with fountain codes. In \emph{2019 27th Signal Processing and Communications Applications Conference (SIU)} (pp. 1–4). IEEE.\par\vspace{3pt}\noindent [CT4] \textbf{Yildiz, H. U.} (2018, November). Minimum delay and maximum lifetime trade-off in underwater sensor networks. In \emph{2018 National Conference on Electrical, Electronics and Biomedical Engineering (ELECO)} (pp. 80–83). EMO.\par\vspace{3pt}\noindent [CT3] Karakurt, Y., \textbf{Yildiz, H. U.}, \& Tavli, B. (2018, May). The impact of mitigation of eavesdropping on wireless sensor networks lifetime. In \emph{2018 26th Signal Processing and Communications Applications Conference (SIU)} (pp. 1–4). IEEE.\par\vspace{3pt}\noindent [CT2] \textbf{Yildiz, H. U.} (2018, May). The impact of data fragmentation on network lifetime in underwater acoustic sensor networks. In \emph{2018 26th Signal Processing and Communications Applications Conference (SIU)} (pp. 1–4). IEEE.\par\vspace{3pt}\noindent [CT1] \textbf{Yildiz, H. U.}, Tavli, B., \& Kahjogh, B. O. (2017, May). Assessment of wireless sensor network lifetime reduction due to elimination of critical node sets. In \emph{2017 25th Signal Processing and Communications Applications Conference (SIU)} (pp. 1–4). IEEE.\par
%%%%%%%%%%%%%%%%%%%%%%%%%%%%%%%%%%%%%%%%%
%-----------CITATIONS-------------------
%%%%%%%%%%%%%%%%%%%%%%%%%%%%%%%%%%%%%%%%%
\clearpage
\vspace{-10pt}
\section{\faChartBar~Citations}
\justify
\noindent $\bullet$ Total \href{https://scholar.google.com/citations?user=nQwHS1gAAAAJ}{\underline{Google Scholar}} Citations: 1,058,
\noindent $\bullet$ \textbf{h-index:} 16,
\noindent $\bullet$ \textbf{i10-index:} 24

%%%%%%%%%%%%%%%%%%%%%%%%%%%%%%%%%%%%%%%%%
%-----------SERVICES-------------------
%%%%%%%%%%%%%%%%%%%%%%%%%%%%%%%%%%%%%%%%%
\vspace{-10pt}
\section{\faSuitcase~Professional Activities and Service}
\justify

\subsection*{(1) \underline{Technical Program Committee Member:}}
\begin{itemize}[leftmargin=*]
\setlength\itemsep{-0.4em}
\item \noindent IEEE International Conference on Communications (ICC 2018–2022),
\item \noindent IEEE Wireless Communications and Networking Conference (WCNC 2019, 2021–2025),
\item \noindent International Balkan Conference on Communications and Networking (BalkanCom 2023–2025),
\item \noindent International Conference on Innovation and Intelligence for Information, Computation, and Technology (3ICT 2019–2020),
\item \noindent International Conference on Network and Service Management (CNSM 2020–2021),
\item \noindent IEEE Middle East and North Africa Communications Conference (MENACOMM 2019),
\item \noindent International Conference on Underwater Networks \& Systems (WUWNet 2019–2024),
\item \noindent IEEE Vehicular Technology Conference (VTC 2018–Spring),
\item \noindent IEEE INFOCOM Workshop on Wireless Communications and Networking in Extreme Environments (2018),
\item \noindent IEEE SENSORS (2017).
\end{itemize}

\vspace{-20pt}

\subsection*{(2) \underline{Reviewer:}}
$\bullet$ Ad Hoc Networks,
$\bullet$ Computer Networks,
$\bullet$ IEEE Access,
$\bullet$ IEEE Communications Letters,
$\bullet$ IEEE Communications Surveys and Tutorials,
$\bullet$ IEEE International Conference on Communications (ICC),
$\bullet$ IEEE Internet of Things Journal,
$\bullet$ IEEE WiMob 2016,
$\bullet$ IEEE PIMRC 2017,
$\bullet$ IEEE Sensors Journal,
$\bullet$ IEEE Systems Journal,
$\bullet$ IEEE Transactions on Cybernetics,
$\bullet$ IEEE Transactions on Industrial Electronics,
$\bullet$ IEEE Transactions on Industrial Informatics,
$\bullet$ IEEE Transactions on Mobile Computing,
$\bullet$ IEEE WCNC 2018–2024,
$\bullet$ IEEE Wireless Communications Letters,
$\bullet$ Internet of Things (Guest Editor),

\vspace{-10pt}
\subsection*{(3) \underline{Session Chair:}}

\begin{itemize}[leftmargin=*]
\setlength\itemsep{-0.4em}
\item \noindent ``PHY-II: Physical Layer Communications-II'' -- IEEE BlackSeaCom 2021,
\item \noindent ``Underwater Networking'' -- ACM WUWNet 2019,
\item \noindent ``Communication Networks II'' -- SIU 2018,
\item \noindent ``Sensor Network, Method and Evaluation'' -- IEEE SENSORS 2016,
\item \noindent ``Management of Emerging Networks'' -- IEEE GLOBECOM 2014 Workshop.
\end{itemize}

%\vspace{-20pt}
%\subsection*{(4) \underline{Editor:}}
%\begin{itemize}[leftmargin=*]
%\setlength\itemsep{-0.4em}
%\item \noindent Computer Standards \& Interfaces, Special Issue on ``Resilience and Security in Next-Generation Wireless Sensor Networks'' (Elsevier, 2025).
%\end{itemize}
%%%%%%%%%%%%%%%%%%%%%%%%%%%%%%%%%%%%%%%%%
%-----------COURSES TAUGHT----------------
%%%%%%%%%%%%%%%%%%%%%%%%%%%%%%%%%%%%%%%%%
\vspace{-15pt}
\section{\faChalkboardTeacher~Courses Taught}
\justify
\subsection*{(1) \underline{TED University:}}
\begin{itemize}[leftmargin=*]
\setlength\itemsep{-0.4em}
\item \noindent EE 205 -- Software Tools for Electrical Engineering (Fall'16),
\item \noindent EE 304 -- Probability and Random Variables (Spring'17 – Present),
\item \noindent EE 311 -- Signals and Systems (Fall'16 – Present),
\item \noindent EE 312 -- Communication Systems I (Spring'17, Spring'18),
\item \noindent EE 413 -- Communication Systems II (Fall'16 – Present),
\item \noindent EE 462 -- Power System Analysis (Spring'17),
\item \noindent EE 512 -- Optimization for Communication Networks (Spring'26).
\end{itemize}

%%%%%%%%%%%%%%%%%%%%%%%%%%%%%%%%%%%%%%%%%
%-----------TALKS-------------------
%%%%%%%%%%%%%%%%%%%%%%%%%%%%%%%%%%%%%%%%%
\vspace{-15pt}
\section{\faUsers~Seminars and Invited Talks}
\justify

\noindent \textbf{Yildiz, H. U.} (2019, October 24). \emph{Prolonging the lifetime of underwater sensor networks under sinkhole attacks} [Invited talk]. Georgia Institute of Technology, Atlanta, GA, USA. Hosted by I. F. Akyildiz.
\vspace{5pt}

\noindent \textbf{Yildiz, H. U.} (2019, April 25). \emph{Maximization of underwater sensor networks lifetime via fountain codes} [Conference presentation]. SIU 2019 — Special Session on Next Generation Communication Techniques, Sivas, Türkiye.
\vspace{5pt}

\noindent \textbf{Yildiz, H. U.} (2019, February 19). \emph{Maximization of underwater sensor networks lifetime via fountain codes} [Invited talk]. ASELSAN Information and Communication Technologies Workshop, Ankara, Türkiye.
\vspace{5pt}

\noindent \textbf{Yildiz, H. U.} (2016, November 14). \emph{Optimal transmission power level sets for lifetime maximization in wireless sensor networks} [Invited talk]. ASELSAN Information and Communication Technologies Workshop, Ankara, Türkiye.
\vspace{5pt}

\noindent \textbf{Yildiz, H. U.} (2014, December 3). \emph{The impact of random power assignment in handshaking on wireless sensor network lifetime} [Invited talk]. University of Texas at Dallas, Dallas, TX, USA. Hosted by M. Torlak.
\vspace{5pt}

\noindent \textbf{Yildiz, H. U.} (2014, February 11). \emph{Transmission power control for link-level handshaking in wireless sensor networks} [Invited talk]. University of Rochester, Rochester, NY, USA. Hosted by W. Heinzelman.
\vspace{5pt}

\noindent \textbf{Yildiz, H. U.} (2014, February 7). \emph{Transmission power control for link-level handshaking in wireless sensor networks} [Invited talk]. University at Buffalo, Buffalo, NY, USA. Hosted by T. Melodia.
\vspace{5pt}

\noindent \textbf{Yildiz, H. U.} (2014, February 5). \emph{A seminar on wireless sensor networks, network optimization, and mathematical programming} [Invited seminar]. Carleton University, Ottawa, ON, Canada. Hosted by H. Yanikomeroglu.
\vspace{5pt}

\noindent \textbf{Yildiz, H. U.} (2014, January 23). \emph{Transmission power control for link-level handshaking in wireless sensor networks} [Invited talk]. University of Southern California, Los Angeles, CA, USA. Hosted by B. Krishnamachari.
\vspace{5pt}

\noindent \textbf{Yildiz, H. U.} (2014, January 22). \emph{Transmission power control for link-level handshaking in wireless sensor networks} [Invited talk]. University of California, Irvine, CA, USA. Hosted by E. Ayanoglu.

%%%%%%%%%%%%%%%%%%%%%%%%%%%%%%%%%%%%%%%%%
%-------THESES SUPERVISED-*--------------
%%%%%%%%%%%%%%%%%%%%%%%%%%%%%%%%%%%%%%%%%
\vspace{-10pt}
\section{\faUserGraduate~Theses Supervised}
\justify

\noindent [TH8] Geleri, S. (in progress). \textit{Joint optimization of node deployment and k-connectivity in underwater acoustic sensor networks for seismic monitoring} (Master's thesis). TED University, Ankara, Türkiye.

\vspace{5pt}

\noindent [TH7] Tantur Karagul, C. (2025). \textit{Network lifetime optimization in underwater wireless sensor networks with variable k-connectivity} (Doctoral dissertation). TOBB University of Economics and Technology, Ankara, Türkiye.

\vspace{5pt}

\noindent [TH6] Cobanlar, M. (2022). \textit{Analysis of the trade-off between network lifetime and k-connectivity in wireless sensor networks} (Doctoral dissertation). TOBB University of Economics and Technology, Ankara, Türkiye.

\vspace{5pt}

\noindent [TH5] Un, B. E. (2021). \textit{Design of a novel optimization framework for analyzing the impact of critical nodes on the network lifetime of underwater acoustic sensor networks using private key cryptography} (Master's thesis). TOBB University of Economics and Technology, Ankara, Türkiye.

\vspace{5pt}

\noindent [TH4] Aydin, C. (2021). \textit{Network lifetime maximization in underwater wireless sensor networks based on the number of depleted nodes} (Master's thesis). TOBB University of Economics and Technology, Ankara, Türkiye.

\vspace{5pt}

\noindent [TH3] Ozmen, A. (2021). \textit{Modeling the trade-off between eavesdropping and network lifetime using a mixed-integer programming approach in underwater acoustic sensor networks} (Master's thesis). TOBB University of Economics and Technology, Ankara, Türkiye.

\vspace{5pt}

\noindent [TH2] Karakurt, Y. (2018). \textit{Eavesdropping potential in wireless sensor networks and modeling and analysis of network lifetime recovery} (Master's thesis). TOBB University of Economics and Technology, Ankara, Türkiye.

\vspace{5pt}

\noindent [TH1] Tantur, C. (2017). \textit{Optimal transmission power level sets for lifetime maximization in wireless sensor networks} (Master's thesis). TOBB University of Economics and Technology, Ankara, Türkiye.


%%%%%%%%%%%%%%%%%%%%%%%%%%%%%%%%%%%%%%%%%
%-----------LANGUAGES----------------
%%%%%%%%%%%%%%%%%%%%%%%%%%%%%%%%%%%%%%%%%
\vspace{-10pt}
\section{\faLanguage~Languages}
\justify

\noindent $\bullet$ \textbf{English:} Business and academic proficiency.
\noindent $\bullet$ \textbf{Turkish:} Native.

%%%%%%%%%%%%%%%%%%%%%%%%%%%%%%%%%%%%%%%%%
%-----------ORGANIZATIONS----------------
%%%%%%%%%%%%%%%%%%%%%%%%%%%%%%%%%%%%%%%%%
\vspace{-10pt}
\section{\faNetworkWired~Professional Memberships}
\justify

\begin{itemize}[leftmargin=*]
\setlength\itemsep{-0.5em}
\item \noindent IEEE Senior Member (2021 – present), IEEE Member (2016 – 2021), IEEE Student Member (2013 – 2016), IEEE Oceanic Engineering Society (2021), IEEE Communications Society (2015 – 2018).
\end{itemize}

%%%%%%%%%%%%%%%%%%%%%%%%%%%%%%%%%%%%%%%%%
%-----------BIOGRAPHY--------------------
%%%%%%%%%%%%%%%%%%%%%%%%%%%%%%%%%%%%%%%%%
%\vspace{-10pt}
%\section{\faEnvira~Biographical Sketch}
%\justify

%Huseyin Ugur Yildiz received his B.Sc. degree from Bilkent University, Ankara, Turkey, in 2009; M.Sc. and Ph.D. degrees from TOBB University of Economics and Technology, Ankara, in 2013 and 2016, all in electrical and electronics engineering, respectively. He worked for Turk Telekom (in 2010-2015) and Turkish Aerospace Industries (in 2015-2016), Ankara, where he designed the network architectures of Safer Internet Service and the ANKA unmanned aerial vehicle. He is currently an associate professor of the electrical and electronics engineering department at TED University, Ankara, Turkey. His research focuses on modeling and solving mathematical optimization problems for wireless ad hoc \& sensor networks, underwater acoustic networks, and smart grids. He is a senior member of the IEEE.

%%%%%%%%%%%%%%%%%%%%%%%%%%%%%%%%%%%%%%%%%
%-----------REFERENCES----------------
%%%%%%%%%%%%%%%%%%%%%%%%%%%%%%%%%%%%%%%%%
\begin{comment}
\vspace{-10pt}
\section{\faAddressCard[regular]~References}
\justify
\begin{itemize}[leftmargin=*]
\setlength\itemsep{-0.5em}
\item \noindent 
\textbf{Bulent Tavli} $\vert$
Professor of Electrical and Electronics Engineering $\vert$
TOBB University of Economics and Technology $\vert$
\faMapMarker* Sogutozu Caddesi No:43, 06560, Ankara, Turkey $\vert$
\faPhone +90 (312) 292 40 74 $\vert$
\faEnvelope~ \href{mailto:btavli@etu.edu.tr}{\underline{btavli@etu.edu.tr}}.

\item \noindent 
\textbf{Vehbi Cagri Gungor} $\vert$
Professor of Computer Engineering $\vert$
Abdullah Gul University $\vert$
\faMapMarker* Sumer Kampusu, 38080 Kayseri, Turkey $\vert$
\faPhone +90 (352) 224 88 00  $\vert$
\faEnvelope~ \href{mailto:cagri.gungor@agu.edu.tr}{\underline{cagri.gungor@agu.edu.tr}}.

\item \noindent
\textbf{Halim Yanikomeroglu} $\vert$
Professor of Systems and Computer Engineering $\vert$
Carleton University $\vert$
\faMapMarker* 1125 Colonel By Dr, Ottawa, ON, Canada $\vert$
\faPhone +1 (613) 520-5734 $\vert$
\faEnvelope~\href{mailto:halim@sce.carleton.ca}{\underline{halim@sce.carleton.ca}}.
\end{itemize}
\end{comment}
%\vspace{20pt}
%\noindent Last updated on \today.

\end{document}
